\documentclass[a4paper,landscape,11pt]{article}

\usepackage{geometry}
\geometry{a4paper,landscape,margin=1cm,top=1cm,bottom=1cm}

\usepackage{amsmath,amssymb}
\usepackage{xcolor}
\usepackage{tikz}
\usetikzlibrary{positioning,arrows.meta,calc,fit,backgrounds,shapes.geometric}

\usepackage[extrafootnotefeatures]{xepersian}

\settextfont[
Scale=1.0,
Extension=.ttf,
Path=/Users/macbook/develope/thesis/thesis_latex/first_edition/styles/fonts/,
BoldFont=XB NiloofarBd,
ItalicFont=XB NiloofarIt,
BoldItalicFont=XB NiloofarBdIt
]{XB Niloofar}

\setdigitfont[
Scale=1.0,
Extension=.ttf,
Path=/Users/macbook/develope/thesis/thesis_latex/first_edition/styles/fonts/,
BoldFont=XB NiloofarBd
]{XB Niloofar}

\pagestyle{empty}
\setlength{\parindent}{0pt}

% ---------- palette ----------
\definecolor{c1}{HTML}{2C5F8A}
\definecolor{c2}{HTML}{A8621B}
\definecolor{c3}{HTML}{2E6B4F}
\definecolor{c4}{HTML}{6B3D8A}
\definecolor{c5}{HTML}{555555}
\definecolor{cred}{HTML}{B02020}
\definecolor{cwarn}{HTML}{B5761A}
\definecolor{cok}{HTML}{2E6B4F}
\definecolor{cgy}{HTML}{6A6A6A}
\definecolor{inkk}{HTML}{1A1A1A}

\newcommand{\bul}{\ensuremath{\bullet}~}
\newcommand{\key}[1]{\textbf{#1}}
\newcommand{\done}{{\color{cok}\bfseries [نوشته‌شده]}}
\newcommand{\wip}{{\color{cwarn}\bfseries [در دست تکمیل]}}
\newcommand{\todo}{{\color{cred}\bfseries [نوشته‌نشده]}}

\tikzset{
  every node/.style={inner sep=5pt, execute at begin node=\setRTL},
  hdr/.style={rounded corners=3pt, draw=#1, line width=1pt, fill=#1!10,
              text width=\colw, align=right, anchor=north, inner sep=6pt},
  arw/.style={-{Stealth[length=3mm,width=2mm]}, line width=0.9pt, draw=black!55},
  arwl/.style={-{Stealth[length=3mm,width=2mm]}, line width=0.9pt, draw=black!35, rounded corners=6pt},
}

\newlength{\colw}

\newcommand{\pagetitle}[2]{%
  \begin{center}
  {\LARGE\bfseries\color{#1} #2}
  \end{center}
  \vspace{-2mm}
}

\begin{document}

% =====================================================================
% PAGE 1
% =====================================================================
\begin{center}
{\LARGE\bfseries نقشه‌ی محتوایی پایان‌نامه}\\[3pt]
{\large آنچه در هر فصل و هر بخش پوشش می‌دهم}
\end{center}
\vspace{1mm}

\begin{center}
\begin{minipage}{25.9cm}
\footnotesize
{\bfseries عنوان:} بخش‌بندی خودکار چنداندامی شکم در تصاویر سی‌تی‌اسکن با تمرکز بر بهبود کیفیت بر روی تصاویر بدون تزریق ماده‌ی حاجب \hfill
{\bfseries نگارش:} سعیده دانائی \quad\textbar\quad {\bfseries استاد راهنما:} دکتر حمیدرضا ربیعی \quad\textbar\quad {\bfseries تاریخ:} شهریور ۱۴۰۵\\[2pt]
\rule{\linewidth}{0.4pt}\\[2pt]
این سند، دامنه‌ی محتوایی هر بخش پایان‌نامه را در یک نگاه خلاصه می‌کند: در هر بخش چه چیزی پوشش داده می‌شود، هر بخش به کدام ادعا یا نتیجه می‌رسد، و وضعیت فعلی نگارش آن چیست. علامت‌ها: \done\ \wip\ \todo
\end{minipage}
\end{center}
\vspace{2mm}

\setlength{\colw}{4.5cm}
\begin{center}
\begin{tikzpicture}[font=\scriptsize]

\node[hdr=c1, minimum height=5.2cm] (a1) at (10.2,0)
 {{\small\bfseries\color{c1} فصل ۱ \raisebox{1pt}{---} مقدمه}\\ \done\\[3pt]
  {\bfseries پرسشی که فصل پاسخ می‌دهد: چرا این مسئله؟}\\[3pt]
  \bul انگیزه‌ی بالینی: فاز بدون‌کنتراست در تروما، غربالگری و نارسایی کلیوی تنها تصویر در دسترس است، اما دقت بخش‌بندی روی آن نظام‌مند پایین‌تر است.\\[2pt]
  \bul تجزیه‌ی مسئله به چهار زیرمسئله، که ستون‌بندی سه فصل بعدی را تعیین می‌کند.};

\node[hdr=c2, minimum height=5.2cm] (a2) at (5.1,0)
 {{\small\bfseries\color{c2} فصل ۲ \raisebox{1pt}{---} کارهای پیشین}\\ \done\\[3pt]
  {\bfseries پرسشی که فصل پاسخ می‌دهد: ادبیات کجا خالی است؟}\\[3pt]
  \bul مرور پنج حوزه: یک حوزه مسئله‌ی چتر را تعریف می‌کند و چهار حوزه‌ی دیگر با چهار زیرمسئله متناظرند.\\[2pt]
  \bul خروجی فصل، چهار جای خالی مشخص است که فصل ۳ به همان ترتیب به آن‌ها می‌پردازد.};

\node[hdr=c3, minimum height=5.2cm] (a3) at (0,0)
 {{\small\bfseries\color{c3} فصل ۳ \raisebox{1pt}{---} روش پیشنهادی}\\ \done\\[3pt]
  {\bfseries پرسشی که فصل پاسخ می‌دهد: چه چیزی ساخته شده است؟}\\[3pt]
  \bul سه مؤلفه که یک زنجیره‌ی وابسته می‌سازند: تشخیص فاز $\leftarrow$ ثبت تصویر بین‌فازی $\leftarrow$ سنتز.\\[2pt]
  \bul به‌علاوه‌ی پروتکل ارزیابی، که به‌عنوان بخشی از سهم روش‌شناختی پایان‌نامه ارائه می‌شود.};

\node[hdr=c4, minimum height=5.2cm] (a4) at (-5.1,0)
 {{\small\bfseries\color{c4} فصل ۴ \raisebox{1pt}{---} آزمایش‌ها}\\ \wip\\[3pt]
  {\bfseries پرسشی که فصل پاسخ می‌دهد: شواهد چه می‌گویند؟}\\[3pt]
  \bul چهار بخش نوشته‌شده: مرور مجموعه‌داده‌ها (۴-۱)، ثبت تصویر (۴-۲)، تشخیص فاز (۴-۳)، و سنتز \lr{PLC-Net} (۴-۴).\\[2pt]
  \bul گزارش شامل ادعاهای تأییدشده و نیز ادعاهایی است که شواهد آن‌ها را رد کرده‌اند.};

\node[hdr=c5, minimum height=5.2cm] (a5) at (-10.2,0)
 {{\small\bfseries\color{c5} فصل ۵ \raisebox{1pt}{---} نتیجه‌گیری}\\ \todo\\[3pt]
  {\bfseries پرسشی که فصل پاسخ می‌دهد: چه به دست آمد و چه باقی است؟}\\[3pt]
  \bul جمع‌بندی یافته‌ها، تصریح محدودیت‌ها (به‌ویژه دامنه‌ی اعتبار ناشی از سوگیری داده) و کارهای آتی.\\[3pt]
  \bul پس از تکمیل بخش‌های تجربی فصل ۴ نوشته می‌شود.};

\draw[arw] (a1.west) -- (a2.east);
\draw[arw] (a2.west) -- (a3.east);
\draw[arw] (a3.west) -- (a4.east);
\draw[arw] (a4.west) -- (a5.east);

\node[draw=black!35, rounded corners=3pt, fill=black!3, text width=25.6cm,
      align=right, anchor=north, inner sep=8pt] (band) at (0,-6.0)
 {{\small\bfseries ستون‌بندی مشترک}\\[2pt]
  چهار زیرمسئله با ترتیبی ثابت در فصل‌های ۲ تا ۴ تکرار می‌شوند، به‌گونه‌ای که مسیر «پیشینه، روش، آزمایش» را بتوان برای هر زیرمسئله به‌صورت مستقل دنبال کرد.};

\setlength{\colw}{5.6cm}
\node[hdr=cgy, minimum height=2.6cm] (b1) at (9.45,-8.0)
 {{\bfseries زیرمسئله ۱: مرور مجموعه‌داده‌ها}\\[2pt]
  پیشینه ۲-۳ \ $\leftarrow$\ روش ۳-۲ \ $\leftarrow$\ آزمایش ۴-۱\\[2pt]
  \done\ \ نتایج آن مستقل نیز منتشر شده است.};
\node[hdr=cgy, minimum height=2.6cm] (b2) at (3.15,-8.0)
 {{\bfseries زیرمسئله ۲: تشخیص فاز کنتراست}\\[2pt]
  پیشینه ۲-۴ \ $\leftarrow$\ روش ۳-۴ \ $\leftarrow$\ آزمایش ۴-۳\\[2pt]
  عمداً فشرده؛ نقش اصلی‌اش ابزار سنجش فصل ۴ است};
\node[hdr=cgy, minimum height=2.6cm] (b3) at (-3.15,-8.0)
 {{\bfseries زیرمسئله ۳: ثبت تصویر بین‌فازی}\\[2pt]
  پیشینه ۲-۵ \ $\leftarrow$\ روش ۳-۵ \ $\leftarrow$\ آزمایش ۴-۲\\[2pt]
  \done\ \ کامل‌ترین زیرمسئله‌ی پایان‌نامه.};
\node[hdr=cgy, minimum height=2.6cm] (b4) at (-9.45,-8.0)
 {{\bfseries زیرمسئله ۴: سنتز \lr{NC\,$\rightarrow$\,CE}}\\[2pt]
  پیشینه ۲-۶ \ $\leftarrow$\ روش ۳-۶ \ $\leftarrow$\ آزمایش ۴-۴\\[2pt]
  \lr{PLC-Net} (نسخه‌ی دوبعدی) \done\ \ مفصل‌ترین بخش تجربی};

\end{tikzpicture}
\end{center}

\vspace{-1mm}
\begin{center}
\fbox{\parbox{25.6cm}{\small\bfseries ادعای کلی پایان‌نامه:\normalfont\ شکاف دقت بخش‌بندی چندعضوی میان فاز بدون‌کنتراست و کنتراست‌دار، بدون بازطراحی مدل بخش‌بندی و تنها با مداخله در \textbf{سطح داده و پیش‌پردازش}، تا اندازه‌ی قابل‌سنجشی کاهش‌پذیر است.}}
\end{center}

\newpage
% =====================================================================
% PAGE 2 : CHAPTER 1
% =====================================================================
\pagetitle{c1}{فصل ۱ \raisebox{1pt}{---} مقدمه: پوشش محتوایی هر بخش}

\setlength{\colw}{7.7cm}
\begin{center}
\begin{tikzpicture}[font=\footnotesize]

\node[hdr=c1, minimum height=7.2cm] (n1) at (8.55,0)
 {{\small\bfseries\color{c1} ۱-۱ اهمیت مسئله} \done\\[4pt]
  در این بخش زمینه‌ی بالینی مسئله را می‌سازم:\\[3pt]
  \bul ساختار چندفازی سی‌تی شکمی و نقش الگوی افزایش کنتراست و شست‌وشو در تشخیص ضایعات کانونی کبد، به‌ویژه \lr{HCC}.\\[3pt]
  \bul استدلال اصلی بخش: فاز بدون‌کنتراست کاربردی صرفاً تکمیلی ندارد؛ در تروما، غربالگری فرصت‌طلبانه و بیماران با نارسایی کلیوی یا سابقه‌ی حساسیت، تنها تصویر در دسترس است.\\[3pt]
  \bul و در عین حال، به دلیل کاهش نسبت کنتراست‌به‌نویز، دقت بخش‌بندی خودکار روی همین فاز به‌طور نظام‌مند و وابسته به اندام پایین‌تر است.};

\node[hdr=c1, minimum height=7.2cm] (n2) at (0,0)
 {{\small\bfseries\color{c1} ۱-۲ تعریف مسئله} \done\\[4pt]
  \bul هدف کلی را تعریف می‌کنم: کاهش شکاف دقت بخش‌بندی چندعضوی میان فاز بدون‌کنتراست و کنتراست‌دار.\\[3pt]
  \bul دو راهبرد مکمل معرفی می‌شوند: هم‌ترازسازی هندسی فازهای موجود از طریق ثبت تصویر، و بازسازی فاز غایب از طریق سنتز.\\[3pt]
  \bul تصریح می‌کنم که هر دو راهبرد بر زیرساختی مشترک تکیه دارند: شناخت محدودیت‌های داده‌ی موجود، و شناسایی صحیح فاز هر بررسی.\\[3pt]
  \bul بر همین مبنا مسئله به چهار زیرمسئله تجزیه می‌شود، که همین چهارتا با همین ترتیب، ستون‌بندی فصل‌های ۲ تا ۴ را تعیین می‌کنند.};

\node[hdr=c1, minimum height=7.2cm] (n3) at (-8.55,0)
 {{\small\bfseries\color{c1} ۱-۳ چالش‌ها} \done\\[4pt]
  چهار چالش فنی مشترک را پوشش می‌دهم و هر یک را به یک تصمیم روشی در فصل ۳ متصل می‌کنم:\\[3pt]
  \bul ناسازگاری فراداده: برچسب فاز در سرآیند \lr{DICOM} ناقص است، پس تشخیص فاز از یک گام جانبی به یک پیش‌نیاز تبدیل می‌شود.\\[3pt]
  \bul تغییر ناهمگن کنتراست: عروق و پارانشیم روشن‌تر می‌شوند و استخوان تقریباً بدون تغییر می‌ماند، و همین فرض بنیادین معیارهای کلاسیک ثبت را نقض می‌کند.\\[3pt]
  \bul کمیابی جفت طبیعی دو فاز برای یک بیمار.\\[3pt]
  \bul ناهمگنی معیار و مجموعه‌داده در ادبیات سنتز، که انتخاب مبنای معتبر را دشوار می‌کند.};

\node[hdr=c1, minimum height=6.4cm] (n4) at (8.55,-7.8)
 {{\small\bfseries\color{c1} ۱-۴ اهداف پژوهش} \done\\[4pt]
  چهار زیرمسئله را به چهار هدف عملیاتی ترجمه می‌کنم و یک هدف پنجم اثر تجمیعی آن‌ها را می‌سنجد:\\[3pt]
  ۱. مستندسازی نظام‌مند سوگیری و کفایت مجموعه‌داده‌های عمومی.\\[2pt]
  ۲. توسعه‌ی مدل تشخیص فاز مستقل از سرآیند.\\[2pt]
  ۳. طراحی پایپ‌لاین ثبت تصویر با محدودسازی دامنه‌ی معیار تراز.\\[2pt]
  ۴. آموزش و ارزیابی مدل‌های سنتز روی جفت‌های هم‌ترازشده.\\[2pt]
  ۵. سنجش اثر مجموع این مداخلات بر شکاف بخش‌بندی.\\[3pt]
  \bul دامنه‌ی کار را نیز صریح می‌کنم: مدل بخش‌بندی بازطراحی نمی‌شود.};

\node[hdr=c1, minimum height=6.4cm] (n5) at (0,-7.8)
 {{\small\bfseries\color{c1} ۱-۵ ساختار پایان‌نامه} \done\\[4pt]
  نقشه‌ی راه فصل‌های بعدی را ارائه می‌کنم:\\[3pt]
  \bul فصل ۲ پیشینه‌ی پنج حوزه و جمع‌بندی جای خالی؛\\[2pt]
  \bul فصل ۳ سه مؤلفه‌ی روش و پروتکل ارزیابی؛\\[2pt]
  \bul فصل ۴ نتایج تجربی هر مؤلفه؛\\[2pt]
  \bul فصل ۵ جمع‌بندی، محدودیت‌ها و کارهای آتی.};

\node[hdr=c1, minimum height=6.4cm] (n6) at (-8.55,-7.8)
 {{\small\bfseries\color{c1} ۱-۶ جمع‌بندی} \done\\[4pt]
  \bul نتیجه‌ی فصل را در یک گزاره جمع می‌کنم: ریشه‌ی افت دقت را نباید تنها در معماری مدل‌ها جست‌وجو کرد.\\[3pt]
  \bul سه عامل مستقل \raisebox{1pt}{---} کاهش کنتراست بافتی، ناسازگاری فراداده‌ی فاز، و کمیابی داده‌ی جفت‌شده \raisebox{1pt}{---} همگی پیش از مرحله‌ی مدل‌سازی عمل می‌کنند.\\[3pt]
  \bul همین گزاره است که مداخله در سطح داده را به‌عنوان راهبرد پایان‌نامه توجیه می‌کند.};

\draw[arw] (n1.west) -- (n2.east);
\draw[arw] (n2.west) -- (n3.east);
\draw[arwl] (n3.south) -- ++(0,-0.3) -- ($(n4.north)+(0,0.3)$) -- (n4.north);
\draw[arw] (n4.west) -- (n5.east);
\draw[arw] (n5.west) -- (n6.east);
\end{tikzpicture}
\end{center}

\newpage
% =====================================================================
% PAGE 3 : CHAPTER 2
% =====================================================================
\pagetitle{c2}{فصل ۲ \raisebox{1pt}{---} کارهای پیشین: پوشش هر بخش و جای خالی حاصل از آن}

\setlength{\colw}{5.6cm}
\begin{center}
\begin{tikzpicture}[font=\scriptsize]

\node[hdr=c2, minimum height=7.2cm] (m1) at (9.45,0)
 {{\small\bfseries\color{c2} ۲-۱ پیش‌زمینه‌ی بالینی} \done\\[4pt]
  \bul فازهای کنتراست و زمان‌بندی آن‌ها را تعریف می‌کنم.\\[3pt]
  \bul توضیح می‌دهم چرا تفسیر رادیولوژیک به بیش از یک فاز نیاز دارد.\\[3pt]
  \bul دامنه‌ی این بخش عمداً محدود است: تنها واژگان و مفاهیمی را می‌سازد که در فصل‌های بعدی به آن‌ها ارجاع داده می‌شود.};

\node[hdr=c2, minimum height=7.2cm] (m2) at (3.15,0)
 {{\small\bfseries\color{c2} ۲-۲ بخش‌بندی چندعضوی و شکاف کیفیت} \done\\[4pt]
  \bul مسئله‌ی چتر پایان‌نامه در این بخش تعریف می‌شود و چهار بخش بعدی زیرمجموعه‌ی آن‌اند.\\[3pt]
  \bul مرور نشان می‌دهد با وجود پیشرفت معماری‌های مبتنی‌بر \lr{nnU-Net}، افت دقت در فاز بدون‌کنتراست باقی است.\\[3pt]
  \bul علت‌های گزارش‌شده در ادبیات: کاهش نسبت کنتراست‌به‌نویز و جابه‌جایی حوزه‌ای وابسته به فاز.};

\node[hdr=c2, minimum height=7.2cm] (m3) at (-3.15,0)
 {{\small\bfseries\color{c2} ۲-۳ مجموعه‌داده‌ها و سوگیری داده} \done\\[4pt]
  \bul چشم‌انداز مجموعه‌داده‌های عمومی و ابزارهای موجود ارزیابی کیفیت داده را پوشش می‌دهم.\\[3pt]
  \bul نتیجه‌ی بخش: هیچ ارزیابی نظام‌مندی از این مجموعه‌داده‌ها به‌عنوان منبع آموزش هوش مصنوعی وجود ندارد، و هیچ منبعی گزارش نکرده چه سهمی از آن‌ها هر دو فاز را در بر می‌گیرند.\\[4pt]
  {\bfseries\color{cred} $\Leftarrow$ جای خالی ۱}};

\node[hdr=c2, minimum height=7.2cm] (m4) at (-9.45,0)
 {{\small\bfseries\color{c2} ۲-۴ تشخیص فاز کنتراست} \done\\[4pt]
  \bul تعریف مسئله، معیارهای ارزیابی رایج و مرور روش‌های موجود.\\[3pt]
  \bul محدودیت‌های ادبیات: اتکا به شبکه‌های سه‌بعدی با نیاز داده‌ی زیاد یا به بخش‌بندی از پیش آموزش‌دیده، و حساسیت به دستگاه و پروتکل تصویربرداری.\\[3pt]
  \bul و ناکارآمدی برچسب‌های سرآیند \lr{DICOM} به‌عنوان مرجع.\\[4pt]
  {\bfseries\color{cred} $\Leftarrow$ جای خالی ۲}};

\node[hdr=c2, minimum height=7.4cm] (m5) at (9.45,-7.8)
 {{\small\bfseries\color{c2} ۲-۵ ثبت تصویر بین‌فازی} \done\\[4pt]
  \bul نشان می‌دهم معیارهای کلاسیک مبتنی‌بر اطلاعات متقابل بر فرض پایداری رابطه‌ی شدت روشنایی در سطح کل تصویر بنا شده‌اند.\\[3pt]
  \bul تزریق ماده‌ی حاجب این فرض را به‌شکل ناهمگن نقض می‌کند.\\[3pt]
  \bul تمایز کلیدی این بخش: روش‌های آگاه از آناتومی موجود، آناتومی را برای وزن‌دهی به کار می‌گیرند و نه برای محدودسازی دامنه‌ی معیار.\\[4pt]
  {\bfseries\color{cred} $\Leftarrow$ جای خالی ۳؛ سهم روشی اصلی پایان‌نامه از همین‌جا شکل می‌گیرد.}};

\node[hdr=c2, minimum height=7.4cm] (m6) at (3.15,-7.8)
 {{\small\bfseries\color{c2} ۲-۶ سنتز \lr{NC\,$\rightarrow$\,CE}} \done\\[4pt]
  \bul مرور از شبکه‌های مولد-متخاصم جفت‌شده تا مدل‌های انتشار اخیر.\\[3pt]
  \bul نتیجه: بازسازی اطلاعات کنتراست ممکن است، اما کیفیت آن به دقت هم‌ترازی هندسی داده‌ی آموزشی وابسته می‌ماند.\\[3pt]
  \bul یک تصور رایج را نیز تصحیح می‌کنم: کد باز در این حوزه به‌کلی غایب نیست.\\[3pt]
  \bul اما هیچ مطالعه‌ای معیاری برای کیفیت هم‌ترازی داده‌ی آموزشی گزارش نکرده است.\\[3pt]
  {\bfseries\color{cred} $\Leftarrow$ جای خالی ۴}};

\node[hdr=c2, minimum height=7.4cm] (m7) at (-6.3,-7.8)
 {{\small\bfseries\color{c2} ۲-۷ جمع‌بندی و جای خالی پژوهش} \done\\[4pt]
  در این بخش، یافته‌های پنج حوزه را کنار هم می‌گذارم و چهار جای خالی مکمل را استخراج می‌کنم:\\[4pt]
  ۱. نیاز به تشخیص قابل‌اعتماد فاز کنتراست، مستقل از برچسب‌های ناقص سرآیند.\\[2pt]
  ۲. نبود پایپ‌لاین ثبت تصویری که دامنه‌ی محاسبه‌ی معیار تراز \raisebox{1pt}{---} و نه صرفاً وزن آن \raisebox{1pt}{---} را به نواحی پایدار محدود کند.\\[2pt]
  ۳. کمیابی جفت‌های هم‌تراز برای آموزش و ارزیابی سنتز.\\[2pt]
  ۴. فقدان شواهد کمّی درباره‌ی اثر این مداخلات بر شکاف بخش‌بندی.\\[4pt]
  \bul ساختار فصل ۳ مستقیماً از همین چهار مورد پیروی می‌کند.};

\draw[arw] (m1.west) -- (m2.east);
\draw[arw] (m2.west) -- (m3.east);
\draw[arw] (m3.west) -- (m4.east);
\draw[arwl] (m4.south) -- ++(0,-0.3) -- ($(m5.north)+(0,0.3)$) -- (m5.north);
\draw[arw] (m5.west) -- (m6.east);
\draw[arw] (m6.west) -- (m7.east);
\end{tikzpicture}
\end{center}

\newpage
% =====================================================================
% PAGE 4 : CHAPTER 3 (A)
% =====================================================================
\pagetitle{c3}{فصل ۳ \raisebox{1pt}{---} روش پیشنهادی (الف): از نمای کلی تا ثبت تصویر}

\setlength{\colw}{7.7cm}
\begin{center}
\begin{tikzpicture}[font=\footnotesize]

\node[hdr=c3, minimum height=7.2cm] (p1) at (8.55,0)
 {{\small\bfseries\color{c3} ۳-۱ نمای کلی} \done\\[4pt]
  \bul منطق زنجیره‌ی مؤلفه‌ها را ارائه می‌کنم: سه مؤلفه یک زنجیره‌ی وابسته تشکیل می‌دهند و نه سه کار موازی.\\[3pt]
  \bul استدلال: آموزش نظارت‌شده‌ی سنتز به جفت‌های هم‌تراز نیاز دارد؛ چنین جفت‌هایی کمیاب‌اند و باید ساخته شوند؛ و ساختن آن‌ها دو پیش‌نیاز دارد: شناسایی فاز و هم‌ترازسازی هندسی.\\[3pt]
  \bul همچنین نشان می‌دهم مدل تشخیص فاز دو نقش دارد: دروازه‌بانی داده، و ابزار اندازه‌گیری در ارزیابی سنتز (بخش ۳-۷).};

\node[hdr=c3, minimum height=7.2cm] (p2) at (0,0)
 {{\small\bfseries\color{c3} ۳-۲ پروتکل مرور نظام‌مند} \done\\[4pt]
  \bul راهبرد جست‌وجو در \lr{Google Scholar}، \lr{PubMed}، \lr{Scopus} و \lr{TCIA} به‌علاوه‌ی مخازن غیرمتعارف.\\[3pt]
  \bul پنج معیار ورود، و ردیابی منشأ هر مجموعه‌داده برای کمّی‌کردن بازاستفاده.\\[3pt]
  \bul چارچوب ارزیابی ریسک سوگیری با هفت مقوله، فرایند دو ارزیاب و سنجش پایایی.\\[3pt]
  \bul دامنه‌ی اعتبار مرور را نیز تصریح می‌کنم: این مرور چشم‌انداز منتشرشده را نمایندگی می‌کند و نه کل منابع موجود.};

\node[hdr=c3, minimum height=7.2cm] (p3) at (-8.55,0)
 {{\small\bfseries\color{c3} ۳-۳ داده و پیش‌پردازش} \done\\[4pt]
  \bul مجموعه‌داده و تقسیم در سطح بیمار با فایل منجمد، که مستقیماً از یافته‌ی نشت داده در بخش ۴-۱ برمی‌آید.\\[3pt]
  \bul استانداردسازی هندسی و شدت روشنایی: بازآرایی به \lr{LPS}، وکسل همسان‌گرد \lr{1.5} میلی‌متر، و پنجره‌ی \lr{HU} واحد در بازه‌ی \lr{$[-200,\,400]$}.\\[3pt]
  \bul استدلال می‌کنم که انتخاب پنجره‌ی واحد یک الزام روش‌شناختی است: مقایسه‌ی دو مدل سنجیده‌شده با دو پنجره‌ی متفاوت معنا ندارد.\\[3pt]
  \bul تمایزی که در ادامه‌ی فصل نقش محوری دارد: نقاب ثبت (دامنه‌ی محاسبه‌ی معیار) در برابر نقاب ارزیابی (دامنه‌ی سنجش).};

\node[hdr=c3, minimum height=7.4cm] (p4) at (8.55,-7.8)
 {{\small\bfseries\color{c3} ۳-۴ تشخیص فاز کنتراست} \done\\[4pt]
  \bul فرض طراحی را صریح می‌کنم: فاز کنتراست، در تعریف بالینی خود، الگوی شدت روشنایی در چند اندام مشخص است.\\[3pt]
  \bul بنابراین بازنمایی ورودی نه حجم خام، بلکه برداری از میانه‌ی \lr{HU} اندام‌هاست؛ انتخاب میانه به‌جای میانگین به دلیل مقاومت در برابر ضایعه و کلسیفیکاسیون است.\\[3pt]
  \bul طبقه‌بند \lr{XGBoost} با دستورالعمل \lr{Segment-and-Classify} و تقسیم در سطح بررسی؛ و آزمایش جانبی کاوش خطی روی رمزگذارهای منجمد \lr{DINOv2} و \lr{MedViT} برای سنجش امکان حذف بخش‌بندی.\\[3pt]
  \bul سه مزیت این طراحی: نیاز به داده‌ی آموزشی اندک، انتقال‌پذیری میان دستگاه‌ها، و تفسیرپذیری سهم هر اندام در تصمیم.};

\node[hdr=c3, minimum height=7.4cm] (p5) at (0,-7.8)
 {{\small\bfseries\color{c3} ۳-۵ الف \raisebox{1pt}{---} ایده‌ی مرکزی ثبت تصویر} \done\\[4pt]
  \bul سهم روشی اصلی پایان‌نامه در این بخش صورت‌بندی می‌شود: معیار تراز نباید در نواحی ناپایدار وزن کمتری بگیرد، بلکه نباید در آن نواحی محاسبه شود.\\[3pt]
  \bul صورت‌بندی صوری: $\sum_{x} w(x)\,s(x)$ در برابر $\sum_{x \in \mathcal{M}} s(x)$.\\[3pt]
  \bul دلیل اهمیت عملی این تفاوت: در فرمول‌بندی وزنی، هیستوگرام مشترک همچنان از وکسل‌های ناپایدار انباشته می‌شود و چندمُدی می‌ماند؛ محدودسازی دامنه، هیستوگرام را نیز محدود می‌کند.\\[3pt]
  \bul ساخت نقاب چهارلایه: استخوان، بافت نرم پایدار، سایر بافت نرم، و ساختارهای حذف‌شده (روده، به‌دلیل حرکت دودی و گاز).};

\node[hdr=c3, minimum height=7.4cm] (p6) at (-8.55,-7.8)
 {{\small\bfseries\color{c3} ۳-۵ ب \raisebox{1pt}{---} پایپ‌لاین چندگذره} \done\\[4pt]
  \bul هم‌ترازسازی مرکز جرم با دروازه‌ی اطمینان و دو ناحیه‌ی مرده، بر این مبنا که یک هم‌ترازی درشتِ نادرست بدتر از نبود آن است.\\[3pt]
  \bul گذر صفر: مقداردهی اولیه‌ی بسته با \lr{SVD} روی مراکز جرم اندام‌ها، همراه با نگهبان دترمینان برای تضمین چرخش سره.\\[3pt]
  \bul گذر یک: پالایش با زیان آناتومیک (دایس فرسوده و مرکز جرم) و بهینه‌ساز بدون مشتق نلدر-مید، که به‌دلیل مشتق‌ناپذیری دایس یک ضرورت است.\\[3pt]
  \bul گذر دو: معیار روی بزرگی گرادیان، با این استدلال که تزریق حاجب سطح روشنایی را جابه‌جا می‌کند اما لبه‌ها را در جای خود نگه می‌دارد.\\[3pt]
  \bul دروازه‌های پذیرش، که بدترین حالت یک گذر ناموفق را از «مخرب» به «بی‌اثر» تبدیل می‌کنند.\\[3pt]
  \bul پالایش نهایی غیرصلب با \lr{B}-اسپلاین؛ تاخوردگی گزارش می‌شود اما به‌عنوان دروازه به کار نمی‌رود.};

\draw[arw] (p1.west) -- (p2.east);
\draw[arw] (p2.west) -- (p3.east);
\draw[arwl] (p3.south) -- ++(0,-0.25) -- ($(p4.north)+(0,0.25)$) -- (p4.north);
\draw[arw] (p4.west) -- (p5.east);
\draw[arw] (p5.west) -- (p6.east);
\end{tikzpicture}
\end{center}

\newpage
% =====================================================================
% PAGE 5 : CHAPTER 3 (B)
% =====================================================================
\pagetitle{c3}{فصل ۳ \raisebox{1pt}{---} روش پیشنهادی (ب): سنتز و پروتکل ارزیابی}

\setlength{\colw}{7.7cm}
\begin{center}
\begin{tikzpicture}[font=\footnotesize]

\node[hdr=c3, minimum height=7.4cm] (q1) at (8.55,0)
 {{\small\bfseries\color{c3} ۳-۶ الف \raisebox{1pt}{---} مولد و شرطی‌سازی} \done\\[4pt]
  \bul صورت‌بندی: ترجمه‌ی تصویر-به-تصویر نظارت‌شده روی جفت‌های هم‌ترازشده‌ی بخش ۳-۵.\\[3pt]
  \bul معماری \lr{U-Net} دوبعدی روی وصله‌های $128 \times 128$، با نرمال‌سازی گروهی؛ دلیل این انتخاب در فصل ۴ به‌صورت کمّی و در ارتباط با مسئله‌ی درز میان وصله‌ها ارائه می‌شود.\\[3pt]
  \bul متمایزکننده‌ی \lr{PatchGAN} شرطی، چون نسخه‌ی غیرشرطی تنها واقع‌نمایی کلی را پاداش می‌دهد و تطابق با ورودی را نمی‌سنجد.\\[3pt]
  \bul تابع زیان ترکیبی، با وزن‌های لایه‌بندی‌شده‌ی سطح اندام؛ روده در آموزش وزن صفر می‌گیرد اما همچنان ارزیابی می‌شود.\\[3pt]
  \bul پیکربندی \LRE{۲٫۵} بعدی برای زمینه‌ی محور $z$، و شرطی‌سازی فاز و سطح با سازوکار \lr{FiLM}.\\[3pt]
  \bul تصریح می‌کنم که شرطی‌سازی سطح یک روش قابل استقرار نیست، بلکه ابزاری برای اندازه‌گیری سقف عملکرد است.};

\node[hdr=c3, minimum height=7.4cm] (q2) at (0,0)
 {{\small\bfseries\color{c3} ۳-۶ ب \raisebox{1pt}{---} عدم‌قطعیت و مدل‌های مبنا} \done\\[4pt]
  \bul استدلال محوری این بخش: رگرسیون آموزش‌دیده با زیان $L_1$ در بهترین حالت میانه‌ی شرطی را پیش‌بینی می‌کند. اگر $y=f(x,Z)$ باشد و $Z$ (دوز، زمان‌بندی بولوس، فیزیولوژی) در فاز بدون‌کنتراست مشاهده‌پذیر نباشد، بخشی از خطا تقلیل‌ناپذیر است.\\[3pt]
  \bul نتیجه: آنچه باید سنجیده شود کاهش خطا نیست، بلکه کالیبراسیون توزیع پیش‌بینی‌شده است.\\[3pt]
  \bul دو خانواده‌ی مدل برای آزمودن این ایده: رگرسیون ناهم‌واریانس به‌عنوان مبنای کم‌هزینه، و مدل انتشار که باید بر آن برتری نشان دهد.\\[3pt]
  \bul انتظار کمّی پیش از اجرای آزمایش ثبت شده است: یک نمونه‌ی صادق از توزیع شرطی، با معیار نقطه‌ای حدود $\sqrt{2}$ بدتر امتیاز می‌گیرد.\\[3pt]
  \bul مدل‌های مبنای ادبیات روی همان تقسیم داده، همان موارد آزمون و همان ابزار امتیازدهی بازآموزی می‌شوند؛ این یک مقایسه‌ی کنترل‌شده است و نه بازتولید نتایج اصلی.};

\node[hdr=c3, minimum height=7.4cm] (q3) at (-8.55,0)
 {{\small\bfseries\color{c3} ۳-۷ الف \raisebox{1pt}{---} نقد معیارهای متعارف} \done\\[4pt]
  این بخش را به‌عنوان بخشی از سهم روش‌شناختی پایان‌نامه ارائه می‌کنم و نه به‌عنوان مقدمه‌چینی. دو دلیل پوشش داده می‌شود:\\[4pt]
  \bul دلیل ساختاری: پنجره‌ی \lr{HU} انتخاب‌شده هوا، ریه و چربی را به صفر و استخوان را به یک اشباع می‌کند؛ معیارهایی که روی کل تصویر میانگین می‌گیرند عمدتاً روی نواحی‌ای میانگین می‌گیرند که همه‌ی مدل‌ها در آن‌ها بی‌خطا هستند.\\[3pt]
  \bul دلیل دوم: \lr{PSNR} و \lr{SSIM} با میانگین شرطی کمینه می‌شوند و بنابراین خروجی تار را بالاتر از بافت صحیح امتیاز می‌دهند، یعنی مدل‌ها را در جهت نادرست مرتب می‌کنند.\\[3pt]
  \bul بر این مبنا، این معیارها گزارش می‌شوند اما ثانویه تلقی می‌شوند.};

\node[hdr=c3, minimum height=7.2cm] (q4) at (8.55,-8.0)
 {{\small\bfseries\color{c3} ۳-۷ ب \raisebox{1pt}{---} معیارهای پیشنهادی} \done\\[4pt]
  \bul معیار اصلی: خطای میانه‌ی \lr{HU} در سطح اندام، که با همان استخراج ویژگی مدل تشخیص فاز محاسبه می‌شود و زنجیره‌ی مؤلفه‌ها را می‌بندد.\\[3pt]
  \bul معیارهای بازیابی سطح: شیب رگرسیون $\beta$ و نسبت واریانس، که یک مولد واقعی را از یک میانگین‌گیر جدا می‌کنند.\\[3pt]
  \bul معیارهای بافت و پیوستگی، که همگی نسبت‌اند و هدفشان مقدار یک است؛ انحراف در هر دو جهت (تاری یا بافت توهمی) شکست تلقی می‌شود.\\[3pt]
  \bul معیارهای کالیبراسیون برای مدل‌های احتمالاتی: \lr{CRPS} و پوشش بازه‌ی پیش‌بینی ۹۰ درصد.\\[3pt]
  \bul \lr{LPIPS} و \lr{FID} تنها برای قابلیت مقایسه با ادبیات گزارش می‌شوند و مبنای هیچ ادعایی نیستند.};

\node[hdr=c3, minimum height=7.2cm] (q5) at (0,-8.0)
 {{\small\bfseries\color{c3} ۳-۷ ج \raisebox{1pt}{---} ارزیابی پایین‌دستی و انضباط آماری} \done\\[4pt]
  \bul تصریح می‌کنم که پرسش عنوان پایان‌نامه تنها با یک ارزیابی پایین‌دستی جداگانه پاسخ داده می‌شود.\\[3pt]
  \bul طرح سه‌بازویی در برابر مرجع (بخش‌بندی روی تصویر کنتراست‌دار واقعی): بازوی بدون‌کنتراست واقعی که خودِ شکاف است، بازوی تصویر تولیدشده، و بازوی شرطی‌شده با سطح واقعی که سقف را می‌دهد.\\[3pt]
  \bul گروه اسکلتی نقش شاهد منفی دارد: افت دایس در این گروه به معنای آن است که آنچه سنجیده شده تفاوت کیفیت کلی تصویر است و نه اثر کنتراست.\\[3pt]
  \bul محدودیت طرح را نیز صریح می‌کنم: مرجع، حاشیه‌نویسی انسانی نیست بلکه خروجی خودکار روی تصویر واقعی است.\\[3pt]
  \bul سه قاعده‌ی آماری در سراسر فصل ۴: کف نویز بذر با دروازه‌ی $2\sigma$، آزمون‌های زوجی در سطح مورد، و تفاضل پیکربندی میان اجراها.};

\node[hdr=c3, minimum height=7.2cm] (q6) at (-8.55,-8.0)
 {{\small\bfseries\color{c3} ۳-۸ جمع‌بندی} \done\\[4pt]
  چهار گزاره‌ی پایانی فصل:\\[4pt]
  ۱. مدل تشخیص فاز، مسئله را از حجم خام به برداری از میانه‌های \lr{HU} در سطح اندام تقلیل می‌دهد و در ادامه به ابزار اندازه‌گیری ارزیابی سنتز تبدیل می‌شود.\\[3pt]
  ۲. پایپ‌لاین ثبت تصویر، به‌جای وزن‌دهی، دامنه‌ی محاسبه‌ی معیار تراز را محدود می‌کند و با مقداردهی اولیه‌ی بسته از کمینه‌های محلی پرهیز می‌کند.\\[3pt]
  ۳. مدل‌های سنتز همگی روی جفت‌های حاصل از این ثبت و روی یک تقسیم منجمد آموزش می‌بینند.\\[3pt]
  ۴. پروتکل ارزیابی خود بخشی از سهم پایان‌نامه است، چون نشان می‌دهد معیارهای متعارف سطح‌پیکسل قادر به رتبه‌بندی مدل‌های این مسئله نیستند.};

\draw[arw] (q1.west) -- (q2.east);
\draw[arw] (q2.west) -- (q3.east);
\draw[arwl] (q3.south) -- ++(0,-0.3) -- ($(q4.north)+(0,0.3)$) -- (q4.north);
\draw[arw] (q4.west) -- (q5.east);
\draw[arw] (q5.west) -- (q6.east);
\end{tikzpicture}
\end{center}

\newpage
% =====================================================================
% PAGE 6 : CHAPTER 4
% =====================================================================
\pagetitle{c4}{فصل ۴ \raisebox{1pt}{---} آزمایش‌ها: یافته‌های گزارش‌شده در هر بخش}

\setlength{\colw}{7.7cm}
\begin{center}
\begin{tikzpicture}[font=\footnotesize]

\node[hdr=c4, minimum height=7.2cm] (r1) at (8.55,0)
 {{\small\bfseries\color{c4} ۴-۱ الف \raisebox{1pt}{---} ترکیب و افزونگی} \done\\[4pt]
  \bul از ۴۵ مجموعه‌داده‌ی واردشده، در ظاهر \LRE{۴۷٬۰۴۹} بررسی سی‌تی فهرست شده است، اما تنها \LRE{۲۳٬۷۴۶} مورد اصیل‌اند؛ یعنی حدود ۵۱ درصد این اکوسیستم بازاستفاده است.\\[3pt]
  \bul ردیابی منشأ نشان می‌دهد این بازاستفاده ساختاری شبکه‌ای دارد: چند مجموعه‌داده‌ی پایه (\lr{LiTS}، \lr{KiTS}، \lr{Pancreas-CT}، \lr{BTCV}) گره‌های مرکزی‌اند و موارد آن‌ها به مجموعه‌های ترکیبی متأخرتر سرریز می‌کند.\\[3pt]
  \bul پیامد روش‌شناختی: خطر نشت داده برای هر پژوهشی که روی مجموعه‌ی ترکیبی آموزش و روی مجموعه‌ی پایه ارزیابی می‌کند.\\[3pt]
  \bul همین ملاحظه، تصمیم فصل ۳ برای تقسیم داده در سطح بیمار را توجیه می‌کند.};

\node[hdr=c4, minimum height=7.2cm] (r2) at (0,0)
 {{\small\bfseries\color{c4} ۴-۱ ب \raisebox{1pt}{---} جغرافیا و ریسک سوگیری} \done\\[4pt]
  \bul با وجود مشارکت ۱۸ کشور، حدود ۷۵ درصد مجموعه‌داده‌ها از ایالات متحده، کانادا و اروپا می‌آیند.\\[3pt]
  \bul مهم‌تر از این تمرکز، غیاب کامل آفریقای زیرصحرا، جنوب آسیا و بخش عمده‌ی خاورمیانه است؛ پیامد بالینی آن را با نمونه‌ی \lr{HCC} در کبد غیرسیروتیک تشریح می‌کنم.\\[3pt]
  \bul ارزیابی ریسک سوگیری روی ۲۲ مجموعه‌داده‌ی دارای صد نمونه یا بیشتر: ریسک بالای نژادی ۱۶ از ۲۲ و انتخاب ۱۵ از ۲۲، در برابر برچسب‌گذاری تنها ۵ از ۲۲.\\[3pt]
  \bul تفسیر: ضعف اصلی این اکوسیستم در کیفیت حاشیه‌نویسی نیست \raisebox{1pt}{---} که تمرکز رایج پژوهش‌های کیفیت داده است \raisebox{1pt}{---} بلکه در ترکیب جمعیت و آسیب‌شناسی داده است.};

\node[hdr=c4, minimum height=7.2cm] (r3) at (-8.55,0)
 {{\small\bfseries\color{c4} ۴-۱ ج \raisebox{1pt}{---} پیامد برای این پایان‌نامه} \wip\\[4pt]
  سه قید عملی که این یافته‌ها بر ادامه‌ی کار می‌گذارند:\\[4pt]
  ۱. ارزیابی روی محک‌های ترکیبی بی‌اعتبار است مگر منشأ هر مورد ردیابی شود؛ بنابراین آزمایش‌ها روی یک مجموعه‌داده‌ی اصیل و منفرد انجام می‌شوند.\\[3pt]
  ۲. کمیابی داده‌ی جفت‌شده، ثبت تصویر و سنتز را از دو گزینه‌ی رقیب به دو ضرورت مکمل تبدیل می‌کند.\\[3pt]
  ۳. سوگیری جغرافیایی و جمعیت‌شناختی، دامنه‌ی اعتبار نتایج را محدود می‌کند و این محدودیت در فصل ۵ دوباره مطرح می‌شود.\\[4pt]
  {\color{cwarn}\bfseries در دست تکمیل: سهم مجموعه‌داده‌هایی که هر دو فاز را در بر می‌گیرند باید از شیت استخراج محاسبه و در متن تثبیت شود.}};

\node[hdr=c4, minimum height=7.4cm] (r4) at (8.55,-7.8)
 {{\small\bfseries\color{c4} ۴-۲ الف \raisebox{1pt}{---} قیف داده و اندازه‌ی مسئله} \done\\[4pt]
  \bul از ۲۶۵ بررسی مجموعه‌داده، تنها ۱۶۹ بررسی (\LRE{۶۳٫۸} درصد) سری بدون‌کنتراست دارند؛ همین عدد، و نه کیفیت هم‌ترازی، سقف تعداد جفت‌های آموزشی سنتز را تعیین می‌کند.\\[3pt]
  \bul دروازه‌ی سلامت بخش‌بندی ۷۷۱ حجم از ۷۷۲ را پذیرفت؛ تنها حجم مردود، هر سه شرط دروازه را هم‌زمان نقض می‌کرد.\\[3pt]
  \bul یافته‌ی خلاف انتظار: میانه‌ی ناهم‌ترازی اولیه تنها \LRE{۲٫۷۹} میلی‌متر است و فازها در حدود ۸۷ درصد موارد از پیش هم‌موقعیت‌اند. این یافته با فرض رایج ادبیات ناسازگار است و تفسیر ابلیشن‌ها را تغییر می‌دهد.\\[3pt]
  \bul گزارش با میانه و چارک انجام می‌شود و نه میانگین، چون توزیع خطا دُم‌دار است؛ و معناداری با آزمون علامت زوجی سنجیده می‌شود.};

\node[hdr=c4, minimum height=7.4cm] (r5) at (0,-7.8)
 {{\small\bfseries\color{c4} ۴-۲ ب \raisebox{1pt}{---} ابلیشن‌ها و آزمون ادعای مرکزی} \done\\[4pt]
  \bul نتیجه‌ی اصلی: محدودسازی دامنه‌ی معیار به نقاب اندام، دقیقاً در همان گروه‌های اندامی که فرضیه پیش‌بینی کرده بود \raisebox{1pt}{---} بافت نرم و عروق \raisebox{1pt}{---} بهبود معنادار می‌دهد.\\[3pt]
  \bul و در گروه اسکلتی، که از پیش به‌عنوان شاهد منفی معرفی شده بود، سودی نمی‌دهد؛ همان رفتاری که فرضیه پیش‌بینی کرده بود و ساختار استدلال را قابل ابطال نگه می‌دارد.\\[3pt]
  \bul دو گذر صلب نخست اثرهای مکمل و بر گروه‌های متفاوت دارند.\\[3pt]
  \bul دروازه‌های پذیرش، یک گذر ناموفق را از عامل خرابی به عامل بی‌اثر تبدیل می‌کنند؛ تمایزی که در عمل تعیین‌کننده بود.};

\node[hdr=c4, minimum height=7.4cm] (r6) at (-8.55,-7.8)
 {{\small\bfseries\color{c4} ۴-۲ ج \raisebox{1pt}{---} مقایسه با روش‌های موجود} \done\\[4pt]
  \bul \lr{DEEDS} در هر سه معیار دقت از روش پیشنهادی بهتر عمل می‌کند ($0.886$ در برابر $0.816$) و تفاوت معنادار است.\\[3pt]
  \bul اما \lr{DEEDS} غیرصلب است و در میانه \LRE{۲٫۹} درصد از وکسل‌های درون نقاب دترمینان ژاکوبین نامثبت دارد؛ روش پیشنهادی صلب است و بنا به ساختار نمی‌تواند تاخوردگی تولید کند. برای کاربرد پایین‌دستی، یعنی ساخت جفت‌های آموزشی، نبود تاخوردگی یک الزام است.\\[3pt]
  \bul مقایسه‌ی هم‌تراز، صلب با صلب: $0.816$ در برابر $0.761$ برای نسخه‌ی بدون نقاب.\\[3pt]
  \bul نکته‌ی جانبی: \lr{uniGradICON} با بهترین خطای کلی، در گروه عروقی $5.8$ میلی‌متر خطا دارد در برابر $3.4$ میلی‌متر روش پیشنهادی.\\[3pt]
  \bul دو ادعای طرح‌شده در نسخه‌ی پیشین که شواهد آن‌ها را رد کرد نیز گزارش می‌شوند: بی‌اثر بودن هم‌ترازی مرکز جرم بر دقت نهایی، و بی‌فایده بودن گذر دو در پیکربندی فعلی.\\[3pt]
  \bul محدودیت: ارقام بر ۲۶ بررسی متناظر استوارند و گسترش ارزیابی به کل ۱۵۶ بررسی ثبت‌شده در دستور کار است.};

\draw[arw] (r1.west) -- (r2.east);
\draw[arw] (r2.west) -- (r3.east);
\draw[arwl] (r3.south) -- ++(0,-0.25) -- ($(r4.north)+(0,0.25)$) -- (r4.north);
\draw[arw] (r4.west) -- (r5.east);
\draw[arw] (r5.west) -- (r6.east);
\end{tikzpicture}
\end{center}

\newpage
% =====================================================================
% PAGE 7 : STATUS & PLAN
% =====================================================================
\pagetitle{inkk}{وضعیت نگارش و برنامه‌ی تکمیل}

\setlength{\colw}{11.9cm}
\begin{center}
\begin{tikzpicture}[font=\footnotesize]

\node[hdr=cred, minimum height=8.6cm] (s1) at (6.3,0)
 {{\small\bfseries\color{cred} بخش‌های نوشته‌نشده}\\[4pt]
  \bul \key{نتایج تجربی تشخیص فاز} (زیرمسئله ۲): بخش ۴-۳ بازنویسی شد: مقایسه با سه روش منتشرشده روی ۶۳ حجم آزمون (طبقه‌بند اندامی \LRE{۰٫۹۸۴}) و آزمایش جانبی بازنمایی‌های عمومی (بهترین: \lr{DINOv2} با \LRE{۰٫۸۴۱}). {\color{cred}باقی‌مانده:} ابزار سنتز روی همه‌ی بیماران آموزش دیده است؛ بازآموزی بدون بیماران آزمون سنتز و امتیازدهی دوباره‌ی \lr{phase} و \lr{prob}.\\[3pt]
  \bul \key{جدول مقایسه‌ی سنتز با ادبیات}: بخش ۴-۴ کامل نوشته شد؛ شکاف پایین‌دستی، مقایسه‌ی زوجی با پنج روش منتشرشده، کنترل فاز، کنترل سطح و اعتبار معیارها همگی با عدد پر شده‌اند. {\color{cred}تنها جدول تفصیلی یازده‌روشه} پس از تثبیت پیکربندی نهایی درج می‌شود.\\[3pt]
  \bul \key{فصل ۵ (نتیجه‌گیری)}: پس از تثبیت پیکربندی نهایی نوشته می‌شود.\\[3pt]
  \bul \key{چکیده‌ی فارسی و انگلیسی}: هنوز تدوین نشده است.\\[3pt]
  \bul \key{شکل ۳-۱ و شکل پایپ‌لاین ثبت}: نمودار جریان کلی روش و نمودار انشعاب شاخه‌ی شاهد باید کشیده شوند.};

\node[hdr=cwarn, minimum height=8.6cm] (s2) at (-6.3,0)
 {{\small\bfseries\color{cwarn} مواردی که نیازمند تصمیم یا اجرای شماست}\\[4pt]
  \bul \key{قید ایست‌بازرسی کالیبراسیون}: دو جاروب کنترل‌پذیری روی \lr{best checkpoint} اجرا شده‌اند، در حالی که \lr{PLC-Net} با \lr{last checkpoint} تعریف شده است. این دو یکی نیستند. یا جاروب دوباره اجرا شود، یا در متن تصریح شود که ارقام از \lr{best} همان اجرا می‌آیند.\\[3pt]
  \bul \key{بازوی بخش‌بندی \lr{PLC-Net}}: جدول ۴-۴ اکنون شش بازو دارد و اثر عروقی در ۲۰ از ۲۰ مورد معنادار است؛ اجرای \lr{TotalSegmentator} روی خروجی مدل نهایی یک سطر باقی‌مانده را پر می‌کند.\\[3pt]
  \bul \key{هزینه‌ی محاسباتی ثبت تصویر}: ابزار اندازه‌گیری در پایپ‌لاین هست اما هیچ خروجی‌ای تولید نشده؛ تا اجرای آن، ادعای «دقت مشابه با زمان و انرژی کمتر» نباید مطرح شود.\\[3pt]
  \bul \key{ناسازگاری ارقام سوگیری}: چکیده‌ی مقاله ۱۷ و ۱۶ می‌گوید و جدول اصلی همان مقاله ۱۶ و ۱۵. ارقام پایان‌نامه از شمارش مستقیم جدول آمده‌اند.\\[3pt]
  \bul \key{دوازده سلول \lr{[عدد؟]} در فصل ۲}: مقالاتی که متن کاملشان در دسترس نبوده یا جدول نتایجشان تصویری است.\\[3pt]
  \bul \key{گسترش مجموعه‌ی ارزیابی ثبت تصویر}: از ۲۶ بررسی متناظر فعلی به کل مجموعه‌ی ثبت‌شده.};

\end{tikzpicture}
\end{center}

\end{document}
